\chapter{АНАЛИЗ ПРЕДМЕТНОЙ ОБЛАСТИ И ПОСТАНОВКА ЗАДАЧИ}

\section{Роль расчёта термодинамических свойств воды в инженерных задачах}

Вода и водяной пар являются одной из ключевых рабочих сред в теплоэнергетике и смежных инженерных областях.  
При проектировании, наладке и эксплуатации оборудования требуется многократно определять состояние рабочей среды по ограниченному набору измеряемых параметров.  
К таким задачам относятся расчёты тепловых схем, подбор режимов работы котлов и турбин, оценка потерь в трубопроводах, моделирование дросселирования и теплообмена, а также анализ переходных процессов.  

В инженерной практике расчётные результаты используются не только как численные значения, но и как основа для построения термодинамических диаграмм и последующего принятия решений.  
Визуальные инструменты (например, диаграммы типа $T$--$s$ и $p$--$h$) позволяют оценивать положение точки относительно купола насыщения и быстро выявлять некорректные или опасные режимы.  

С точки зрения программной реализации важны два обстоятельства.  
Во-первых, вычисления должны быть достаточно быстрыми, чтобы поддерживать интерактивные сценарии и пакетные расчёты.  
Во-вторых, алгоритм должен быть устойчивым на границах фаз и в областях резкой нелинейности свойств, поскольку именно там ошибки модели наиболее критичны.  

\section{Стандарты и существующие подходы к расчёту свойств воды и пара}

Расчёт свойств воды и пара исторически выполнялся по табличным данным (паровым таблицам) и диаграммам, в том числе с использованием интерполяции.  
Современные программные средства, как правило, реализуют один из двух подходов: прямое вычисление по аналитическим зависимостям или комбинирование формул с численными методами для обратных маршрутов.  

При инженерных расчётах широко применяются стандартизованные формулировки IAPWS, предназначенные для воды и водяного пара в различных областях состояния.  
Для промышленного применения важен компромисс между точностью и вычислительной эффективностью, а также чётко определённые границы применимости.  

Ключевой особенностью предметной области является разбиение рабочей области на регионы с разными уравнениями состояния и разной «естественной» постановкой задачи.  
Например, в высокоплотной околокритической области (Region~3) удобнее использовать представление через плотность и температуру, что приводит к необходимости обратных зависимостей и/или численных процедур при задании состояния в других параметрах.  
Для практических реализаций это означает, что вычислительный модуль должен включать не только формулы, но и алгоритмы маршрутизации и проверки корректности ввода.  

\section{Обоснование выбора IAPWS-IF97 как основы модели}

В данной работе в качестве базовой физической и нормативной основы используется стандарт IAPWS-IF97 в редакции 2012 года \cite{iapws_if97_2012}.  
Данный стандарт ориентирован на промышленное применение и задаёт набор уравнений и обратных зависимостей для регионов 1--5, линию насыщения (Region~4) и правила топологии области расчёта.  

Использование IAPWS-IF97 обеспечивает согласованность результатов в широкой области давлений и температур, а также совместимость с общепринятыми инженерными расчётами.  
Региональная структура IF97 позволяет строить вычислительные алгоритмы, которые выбирают нужную модель автоматически после классификации состояния.  

При этом стандарт не задаёт единой универсальной аналитической зависимости для всех возможных входных пар параметров на всей области применимости.  
Так, для Region~3 существуют специализированные обратные зависимости (например, для удельного объёма как функции давления и температуры) \cite{iapws_vpt_region3}, а для маршрута $(\rho, T)$ требуется отдельная численная постановка.  
Поэтому в рамках проекта IF97 используется как основа, а отдельные маршруты дополняются численными методами без изменения физической модели.  

\section{Требования к программному комплексу}

Требования к разрабатываемому программному комплексу формулируются исходя из инженерных сценариев использования и особенностей стандарта IF97.  
Комплекс должен обеспечивать единый источник вычислений и одинаковые результаты во всех внешних каналах доставки.  

Функциональные требования:
\begin{itemize}
\item расчёт термодинамического состояния воды и водяного пара в области применимости IAPWS-IF97;  
\item поддержка прямых и обратных маршрутов расчёта по основным входным параметрам ($p$, $T$, $h$, $s$, $\rho$, $x$);  
\item автоматическое определение региона IF97, включая обработку линии насыщения и границы B23;  
\item возврат унифицированного набора выходных свойств (например, $v$, $\rho$, $h$, $s$, $u$, $c_p$, $w$) и номера региона;  
\item проверка диапазонов входных параметров и корректная обработка ошибок ввода и особых случаев.  
\end{itemize}

Нефункциональные требования:
\begin{itemize}
\item вычислительная эффективность для интерактивного использования и пакетных расчётов;  
\item детерминированность результатов и воспроизводимость сборки;  
\item возможность интеграции ядра в настольные приложения и сервисный контур без дублирования предметной логики;  
\item расширяемость архитектуры за счёт разделения ядра и thin-layer адаптеров доставки результата.  
\end{itemize}

Архитектурно проект реализуется по схеме «одно вычислительное ядро и несколько каналов доставки».  
Вычислительное ядро \texttt{if97\_core} реализует расчёт и валидацию.  
Слой DTO \texttt{if97\_app\_api} фиксирует контракты обмена между компонентами.  
Пользовательские приложения и сервис используют ядро через стабильный интерфейс без дублирования предметной логики.  

\section{Постановка вычислительной задачи и поддерживаемые режимы ввода-вывода}

Постановка вычислительной задачи заключается в определении термодинамического состояния воды или водяного пара по заданным входным параметрам и в расчёте полного набора свойств.  
В качестве входа рассматриваются параметры $p$, $T$, $h$, $s$, $\rho$, $x$, которые типичны для инженерных измерений и расчётов.  

В качестве основного результата требуется вычислять следующие величины: удельный объём $v$, плотность $\rho$, энтальпию $h$, энтропию $s$, внутреннюю энергию $u$, изобарную теплоёмкость $c_p$, скорость звука $w$, а также классификацию по регионам IAPWS-IF97.  
Дополнительно требуется обеспечить корректную обработку двухфазной области и линии насыщения (Region~4), поскольку многие инженерные сценарии включают кипение, конденсацию и анализ влажности пара.  

Поддерживаемые режимы ввода-вывода (маршруты расчёта) включают:
\begin{itemize}
\item прямой маршрут $pT$: расчёт состояния по давлению и температуре;  
\item обратные маршруты $pH$ и $pS$: расчёт по давлению и энтальпии или по давлению и энтропии;  
\item маршрут $pX$: расчёт по давлению и степени сухости в двухфазной области;  
\item маршрут $\rho T$: расчёт по плотности и температуре, реализованный как численное расширение (подбор давления методом секущих с последующим вычислением по $pT$).  
\end{itemize}

Во всех режимах должна выполняться валидация входных параметров и контроль области применимости.  
При невозможности расчёта (например, из-за выхода за границы IF97 или некорректной комбинации параметров) система должна возвращать диагностируемую ошибку, а не неопределённые значения.  
